\begin{python0}
from solutions import *; clear()
\end{python0}
\sectionthree{Root finding: An inefficient method}

Another very common problem in Math and Sciences is finding solution(s)
to an equation such as finding solutions to quadratic equations:

\[2 * x * x + 4 * x - 3 = 0\]

This is easy since we have the quadratic equation formula. But this is
highly non-trivial if the equation is more complex such as a degree 5
equation:

\[2 * x * x * x * x * x + 4 * x * x - 3 * x + 10 = 0\]

or even one with trigonometric, logartihmic, and/or exponential
functions!!!

For instance, what if I want to find a solution to

\[3 * x * x - 10 * x * sin(x) - 500 = 0\]

in the range [0, 20]? You don't recall a solution to
such an equation from your math classes right?!?

Since we don't have a formula to use, we just have to
plot the graph and visually see where the function touches the x-axis.
In the same way, we can compute the values of the above expression

\[3 * x * x - 10 * x * sin(x) - 500\]

for many many many many many values in the range [0, 20] and see
which value of x will give a value of

\[3 * x * x - 10 * x * sin(x) - 500\]

that is very close to 0. For instance we can test 1,000,000 values in
the range [0,20]. If we do this by hand, we would go crazy. But
\ldots

We know C++!!!!

All we need to do is to run a for-loop of 1,000,000 values in the range
[0,20] and find the value of x that will get use as close as
possible to the y-axis. Close to the y-axis means

\[|f(x)|\]

is close to zero. In other words, if you're finding the
point closest to the y-axis it means that you need to find the smallest
$|f(x)|$ for the 1,000,000 x values in the range.

